Denk an den Fall. Nullniveau der Lageenergie: der Boden. Zustand 1: das fallende Auto am Anfang, in Ruhe (nur Lageenergie). Zustand 2: kurz vor dem Aufprall (nur kinetische Energie).

\begin{abcliste}
\abc
Sei $h$ die Höhe und $v$ die Geschwindigkeit. Energieerhaltung zwischen den Zuständen 1 und 2:
\begin{align}
 m\,g\,h &= \frac{1}{2}\,m\,v^2
\end{align}
Mit $v = \vO = \v$ (km/h durch 3.6 teilen):
\begin{align}
 h &= \hF \\
   &= \frac{\left(\v\right)^2}{2\cdot\g} \\
   &= \h \approx \result{\hP{0}{2}}
\end{align}

\abc
Die Höhe wächst mit dem Quadrat der Geschwindigkeit: Doppelte Geschwindigkeit braucht die vierfache Höhe.
\begin{align}
 h' &= \hbF = 4\,h \\
   &= \frac{\left(\vb\right)^2}{2\cdot\g} \\
   &= \hb \approx \result{\hbP{0}{2}}
\end{align}
Das sind etwa \nStP{0}{2} Stockwerke eines Gebäudes (mit \SI{3}{\meter} pro Stockwerk).
\end{abcliste}
